\section{Introduction}
\IEEEPARstart{T}{elevision} audiences no longer watch broadcasts in silence. During live debates, political talk shows, and reality competitions, viewers actively turn to second-screen platforms like Twitter to share commentary, memes, and screenshots \cite{SocialTV}. This stream of real-time viewer feedback offers broadcasters far richer signals than traditional audience ratings, providing immediate insight into public controversy, approval, and emotional engagement \cite{memorai}. 

Capturing these affective signals computationally, however, presents severe challenges. Social TV commentary is short, colloquial, and inherently multimodal. Viewers attach screengrabs and reaction images that do not merely mirror the written text; they exaggerate emotional tone, introduce sarcasm, or invert the literal message completely \cite{schifanella2016multimodal, cai2019multimodal}. 

The foundational question that motivated this investigation is: \textit{how much do attached images genuinely contribute to identifying human affective intent in multimodal social network posts?} \cite{poria2017review} During the manual annotation of our dataset, we observed a crucial empirical phenomenon: human annotator agreement increased substantially when annotators were provided with the attached images compared to reading the text alone. Visual context anchored ambiguous social commentary, establishing consensus where text alone yielded divergent interpretations. This observation directly motivates our central research hypothesis: \textit{does this visual grounding similarly carry over to neural models, and does it hold equally for literal and figurative discourse?}

Furthermore, while multilingual initiatives have advanced text-based emotion analysis, computational resources for Italian remain scarce, and existing benchmarks focus exclusively on unimodal text \cite{bianchi2021feel, emit, MultiEmotion}. To bridge this gap, we introduce \textbf{MultiEMiT}, a multimodal dataset of Italian Twitter commentary gathered from television broadcasts. MultiEMiT provides annotations across five pragmatic layers: fine-grained emotions under Plutchik's taxonomy \cite{Plutchik}, irony and sarcasm, offensiveness, target orientation (program container versus broadcast content), and semiotic text-image relationships \cite{MultiM}. Our benchmarking focuses on three primary classification tasks: Emotion Detection, Irony Detection, and Semiotic Relation Detection, with offensiveness and target serving as contextual signals. To expand training coverage for baseline text encoders, we also draw on the external unimodal EMiT (for Emotions) \cite{emit} and IronITA (for Irony) \cite{ironita2018} corpora.

An in-depth statistical analysis of MultiEMiT's multi-layered annotations revealed another compelling pattern: extensive empirical co-occurrences connect different annotation levels. For instance, irony correlates with joy, sarcasm pairs strongly with disgust and offensiveness, and program-directed commentary associates with viewer trust. Rather than treating these dimensions as isolated tasks, we investigated how much information each annotation layer carries about the others. To this end, we propose \textbf{Correlation Prompting}, a strategy that exposes the remaining annotation layers to the model in natural language: trained encoders receive verbalized auxiliary annotations prepended to each post and learn their association with the target labels, while generative models used without task-specific training additionally receive the corpus-level Pearson correlation rules in their instructions. Since gold auxiliary labels are supplied at inference time, this setting serves as an oracle probe of the informativeness of each annotation layer rather than as a deployable system.

Since our goal is to characterize the resource and its possibilities rather than to maximize performance, we frame our study around three core research questions:
\begin{itemize}
    \item \textbf{RQ1 (Visual Context):} Does adding visual context help neural models detect emotion, irony, and semiotic relations, and does its contribution differ across tasks?
    \item \textbf{RQ2 (Correlation Prompting):} Does exposing the auxiliary annotation layers in the model input improve classification across discriminative and generative architectures?
    \item \textbf{RQ3 (Architectural Paradigms):} How do compact discriminative models (BERT-ViT Dual-Encoders with linear classifiers on frozen features) compare with a 12 billion parameter vision--language model (Pixtral 12B) on a shared test set?
\end{itemize}

The main contributions of this work are summarized as follows:
\begin{itemize}
    \item We introduce \textbf{MultiEMiT}, the first multimodal dataset of Italian Social TV commentary, annotated across five pragmatic layers and validated through three-annotator adjudication.
    \item We show that the role of visual context depends strongly on the task: images directly help irony detection by capturing text--image contrast, but help emotion recognition only when combined with auxiliary task labels.
    \item We analyze \textbf{Correlation Prompting} as an exploratory method to examine how much information is shared across annotation layers, testing how auxiliary labels guide both traditional classifiers and generative models.
    \item We compare compact Dual-Encoder baselines with the 12B parameter Pixtral model on a shared test set, highlighting differences in how they handle class imbalance and complex semiotic relations.
\end{itemize}

The rest of this paper is organized as follows. Section~\ref{sec:related_works} situates MultiEMiT within Italian affective computing, Social TV analytics, and multimodal document intent. Section~\ref{sec:dataset_analysis} describes the corpus collection, annotation protocol, and cross-dimensional correlations. Section~\ref{sec:experimental_design} defines our classification tasks, Dual-Encoder architecture, correlation prompting strategy, and Pixtral configurations. Section~\ref{sec:results} presents comparative evaluations and subset analyses across semiotic relations. Section~\ref{sec:conclusion} summarizes our findings and directions for future work.